D'un cartró quadrat de $12$ cm de costat es retallen quatre quadrats iguals, de costat $x$, a
les cantonades. Després es dobleguen les vores cap amunt per fer una capsa sense tapa.

\begin{apartats}

\apartat[1,25]{1}
Comprova que el volum de la capsa és $V(x)=x(12-2x)^2$ i digues entre quins valors pot variar
$x$.

\begin{solucio}
La base de la capsa és un quadrat de costat $12-2x$ i l'alçada és $x$. Per tant,
\[
V(x)=x(12-2x)^2 .
\]
Cal que el costat retallat sigui positiu i que en quedi base: $x>0$ i $12-2x>0$, és a dir
$x\in(0,6)$.
\end{solucio}

\apartat[1,25]{0,75}
Troba el valor de $x$ que fa màxim el volum i calcula aquest volum.

\begin{solucio}
$V(x)=4x(6-x)^2$, i derivant el producte:
\[
V'(x)=4(6-x)^2+4x\cdot2(6-x)(-1)=4(6-x)\left[(6-x)-2x\right]=4(6-x)(6-3x).
\]
S'anul·la en $x=6$ (fora del domini) i en $x=2$.\\
$V(2)=2\cdot8^2=128$: cal retallar quadrats de $\boxed{2}$ cm i el volum màxim és
$\boxed{128}$ cm$^3$.
\end{solucio}

\begin{nomesllarg}
\apartat{0,75}
Justifica que el valor trobat dona un màxim.

\begin{solucio}
A $(0,2)$, $V'>0$ (els dos factors són positius) i a $(2,6)$, $V'<0$ (perquè $6-3x<0$): la
funció creix i després decreix, i per tant en $x=2$ hi ha un \textbf{màxim}.\\
També es pot comprovar amb la derivada segona: $V''(x)=24x-96$ i $V''(2)=-48<0$.
\end{solucio}
\end{nomesllarg}

\end{apartats}
